\documentclass[12pt]{article}

\usepackage[margin=1in]{geometry}
\usepackage{graphicx}
\usepackage[dvipsnames]{xcolor}
\usepackage{enumitem}
\usepackage{parskip}

% ── Fonts and spacing ──
\renewcommand{\familydefault}{\sfdefault}
\pagestyle{empty}

% ── Colors ──
\definecolor{accent}{HTML}{2E5090}
\definecolor{lightgray}{HTML}{F2F2F2}

\begin{document}

% ── Title ──
\begin{center}
  {\LARGE\bfseries\color{accent} Why Your Job Affects How You Should Invest}\\[0.6em]
  {\large Engineer Investor}\\[0.3em]
  {\small\color{gray} \today}
\end{center}

\vspace{0.8em}
\hrule height 0.4pt
\vspace{1.0em}

% ── Section 1 ──
{\large\bfseries\color{accent} The Conventional Wisdom}

Most financial advice tells young people to invest aggressively in stocks.
The reasoning sounds intuitive: you have decades of future paychecks ahead of you,
and that stream of income acts like a giant bond sitting in the background of your
finances. If the market crashes, your salary keeps coming in, giving you plenty of
time to recover. So load up on stocks while you are young.

% ── Section 2 ──
{\large\bfseries\color{accent} The Problem}

That logic breaks down when your income is tied to the stock market. If you work in
tech and a large share of your compensation comes from restricted stock units (RSUs),
your pay rises and falls with equities. The same is true if you work in a cyclical
industry prone to layoffs during downturns, or if your bonus depends on corporate
profits. In these cases, your future earnings are not stable or bond-like at all.
They behave more like stocks, which means you already have enormous, hidden equity
exposure through your career.

% ── Section 3 ──
{\large\bfseries\color{accent} The Fix}

The human capital beta framework solves this with a single number. Beta measures how
closely your future earnings move with the stock market. A beta near zero means your
income is steady regardless of market conditions. A high beta means your income swings
in lockstep with equities. The model splits your human capital into a ``safe'' portion
and a ``risky'' portion, then adjusts your recommended portfolio accordingly. The
higher your beta, the less your career can cushion stock market losses, so the model
tells you to hold a more conservative portfolio.

% ── Section 4 ──
{\large\bfseries\color{accent} What It Means for You}

\begin{itemize}[leftmargin=1.5em, itemsep=0.6em]
  \item \textbf{Government worker} (beta $\approx$ 0): Invest aggressively in your
    portfolio. Your career income is very stable and acts like a large bond, giving
    you capacity to take more risk with your investments.

  \item \textbf{Tech worker with RSUs} (beta $\approx$ 0.6): Moderate your portfolio
    risk. A meaningful share of your compensation already moves with the stock market,
    so piling on more equity exposure in your brokerage account doubles down on the
    same bet.

  \item \textbf{Startup founder} (beta $\approx$ 0.8): Be conservative in your
    portfolio. Your entire career is already a large equity bet. Your savings should
    provide ballast, not amplify that risk.
\end{itemize}

\vfill

% ── Footer ──
\hrule height 0.4pt
\vspace{0.5em}
{\footnotesize\color{gray}
\textbf{Disclaimer:} This document is for educational and research purposes only.
It does not constitute investment advice. Consult a qualified financial advisor before
making investment decisions.\\[0.3em]
\textbf{Source:} \texttt{github.com/engineerinvestor/lifecycle-allocation} \enspace
\textbar\enspace Based on Choi et al., ``Lifecycle Portfolio Choice'' (NBER Working Paper).}

\end{document}
